% Memorandum of Understanding with an Institution (BINDING). RFC-0136 D3, D7.
%
% NOT ISSUABLE BEFORE INCORPORATION, and that is enforced twice:
%   1. `docket counterparty issue <slug> MOU` refuses while the issuing company's `cin` is null in
%      registry/org.yaml (`requires_incorporation: true` in lifecycle.yaml).
%   2. This template reads \dkOrgCin and \dkOrgRegisteredOffice bare. Both are null in org.yaml
%      until incorporation, a null defines nothing, and pdflatex stops on the undefined command
%      (paridyum.cls: a signed document never carries a silently blank field). Never
%      \providecommand either of them.
% Until then the non-binding Letter of Collaboration (ICL) is the document to issue.
%
% Issued to an ORGANISATION COUNTERPARTY (`recipient: organisation`): the party block reads
% \dkCLegalName, \dkCAddress, \dkCOffice about the institution. Whoever signs for it writes their
% name by hand: docket holds no person's data for an institution (RFC-0136 open question 3).
%
% THE NEGOTIATED TERMS ARE AN INPUT, read bare as \dkT... from the `--terms` file the MOU's
% lifecycle entry declares (`domain::terms`). No \dkfill: a registered, hash-verified MOU that is
% a blank form proves nothing about what was agreed (review M5). The institution's legal
% description and the body it acts through are terms too -- the contracting entity's capacity
% matters legally (COUNSEL Q11) -- and every contact is an office or a role mailbox, never a person.
%
% From knowledge docs/2026-11-07-personas-and-paridyum-in/08-paperwork/05.
%
% DRAFT FOR REVIEW BY QUALIFIED COUNSEL (06-counsel-brief.md Q11, Q12, Q13, Q17). Not legal advice.
% Sources: BET-0010, RFC-0127 D4/D7, REQ-0003, STD-0021, RFC-0134 D9/D12, DPDP Act 2023.
%
% The shared governing-law clause says "This Agreement", so this MOU defines itself as the
% "Agreement" in its opening paragraph; that keeps the shared clause correct without a copy.
\documentclass{paridyum}
\begin{document}

\paridyumletterhead
\paridyumdateref{\dkDocDate}{\dkDocRef}
\paridyumverifiable{\dkDocVerify}
\paridyumtitle{Memorandum of Understanding}

This Memorandum of Understanding (``Agreement'') is entered into on the date above between:

\vspace{0.5em}
\begin{tabularx}{\linewidth}{@{}F{3.4cm}X@{}}
Paridyum: & \textbf{\dkOrgName}, a company incorporated under the Companies Act, 2013,
            with Corporate Identity Number \dkOrgCin, whose registered office is at
            \dkOrgRegisteredOffice{} (``Paridyum''); and \\[0.6em]
Institution: & \textbf{\dkCLegalName}, \dkTLegalDescription, whose address is \dkCAddress,
            acting through its \dkTActingThrough{} (the ``Institution''), notices to which are
            addressed to: \dkCOffice.
\end{tabularx}

\section*{1. Purpose}
The parties wish to support students of the Institution who want to start and run businesses.
Paridyum will contribute talks, workshops and written material. The Institution will make room
for them in the calendar of its \dkTActingThrough.

\section*{2. Paridyum's Obligations}
Paridyum will:
\begin{itemize}[leftmargin=*, noitemsep]
    \item deliver the talks and workshops listed in \textbf{Schedule A}, or others agreed in
          writing, at the agreed dates;
    \item provide speakers who describe Paridyum's services accurately and as they exist on the
          date of the session;
    \item provide written material the Institution may reuse with its students, without tracking
          links; and
    \item make its services available to students and student ventures of the Institution only
          on Paridyum's standard published terms, in their own right.
\end{itemize}

\section*{3. The Institution's Obligations}
The Institution will:
\begin{itemize}[leftmargin=*, noitemsep]
    \item provide a calendar slot, and a room or an online meeting, for each session in
          Schedule A;
    \item name a contact for scheduling, set out in Schedule A; and
    \item tell Paridyum promptly if a session cannot go ahead.
\end{itemize}

\section*{4. No Payment}
No money is payable by either party to the other under this Agreement. Paridyum pays no
sponsorship or fee and charges none. Neither party will offer students prizes in credits,
waived charges or service capacity in connection with this Agreement. Each party bears its own
costs.

\section*{5. No Exclusivity}
Nothing in this Agreement prevents either party from working with anyone else, on any terms.

\section*{6. Intellectual Property}
Each party keeps all rights in its own names, logos, materials and other intellectual property.
Paridyum grants the Institution a non-exclusive, royalty-free licence to reproduce and share the
written material it provides under clause 2, unmodified, with the Institution's students and
staff, for non-commercial educational purposes, for as long as this Agreement is in force.
Nothing else is licensed by either party.

\section*{7. Names and Logos}
Neither party may use the other's name, logo, emblem or other marks, or describe the other as a
partner, sponsor or endorser, except exactly as listed in \textbf{Schedule B}. Any other use
needs that party's prior written permission. Either party may withdraw a permission in Schedule
B by email; the other will stop new uses within \dkTLogoWithdrawDays{} days.

\section*{8. Confidentiality}
Information one party marks as confidential, or that a reasonable person would understand to be
confidential, and that it gives the other under this Agreement, will be used by the other only
for this Agreement and not disclosed, except where required by law. This does not apply to
information that is public without a breach of this Agreement, was already lawfully known to
the receiving party, or is independently developed by it. This clause continues for
\dkTConfidentialityYears{} years after this Agreement ends.

\section*{9. Data Protection}
\begin{infobox}
\textbf{No personal data is exchanged under this Agreement.}
\begin{itemize}[leftmargin=*, noitemsep]
    \item Neither party will give the other, and neither will ask for or accept from the other,
          any personal data of students, staff, attendees or registrants --- including names,
          contact details, lists, rosters or r\'esum\'es.
    \item Paridyum will not collect personal data at any session, and will not run sign-ups,
          contests or registrations there. A student who wants to use Paridyum's services does so
          directly, under Paridyum's own terms and privacy notice.
    \item Each party is a separate data fiduciary for any personal data it processes in its own
          operations, and is solely responsible for that processing under the Digital Personal
          Data Protection Act, 2023. Neither party processes personal data on the other's behalf.
    \item The only personal data either party holds about the other is the business contact
          details of the people named in Schedule A, used only to administer this Agreement.
\end{itemize}
\end{infobox}

\section*{10. Students Under Eighteen}
Sessions under this Agreement are for students of the Institution's higher-education
programmes. Nothing under this Agreement will be delivered at a school or directed at persons
under eighteen, unless the parties agree it in writing after Paridyum has obtained legal advice
that it may do so.
% src: RFC-0134 D9, D12; COUNSEL Q7.

\section*{11. Term and Termination}
This Agreement begins on the date above and continues for \dkTTermMonths{} months
unless ended earlier. Either party may end it at any time, without giving a reason, by written
notice (email is enough) to the other's contact in Schedule A, effective
\dkTNoticeDays{} days after it is sent. Clauses 6 (as to rights already exercised),
7, 8 and 9 survive its end.

\section*{12. No Partnership or Agency}
Nothing in this Agreement creates a partnership, joint venture, agency or employment
relationship. Neither party may bind the other.

\section*{13. Entire Agreement and Changes}
This Agreement, with its Schedules, is the whole agreement between the parties on its subject.
It replaces any earlier letter of collaboration between them. It may be changed only in writing
signed for both parties.

\section*{14. Governing Law}
% Governing law. ONE HOME, used by five documents.
%
% The body only -- each document keeps its own `\section*{N. Governing Law}`, because the
% number differs between them and a shared file cannot know which document it is in.
%
% Extracted because it was five copies of one clause, two of which differed only in where the
% line wrapped. A jurisdiction change is one edit rather than five, and there is no fifth copy
% left behind for somebody to find in a year.
%
% Both values come from `registry/org.yaml`, which is the single home for Paridyum's own facts.
This Agreement will be governed by the laws of \dkOrgJurisdictionCountry. Subject to
applicable law, disputes relating to this Agreement will be subject to the appropriate
courts in \textbf{\dkOrgJurisdictionForum}.


\section*{15. Signatures}
Each person signing confirms that they are authorised to sign for the party named above them.

\vspace{1.5em}
\paridyumsignblock

\vspace{1.5em}
\textbf{For \dkCLegalName}

\vspace{0.5em}
\begin{tabularx}{\linewidth}{@{}F{3.4cm}X@{}}
Name:        & \dkblank[6cm] \\[1em]
Designation: & \dkTInstitutionSignatoryOffice \\[1em]
Signature:   & \dkblank[6cm] \\[1em]
Date:        & \dkblank[6cm]
\end{tabularx}

\newpage
\section*{Schedule A --- Sessions and Contacts}
\begin{tabularx}{\linewidth}{@{}F{4.6cm}X@{}}
Sessions planned:            & \dkTSessions \\[0.5em]
Paridyum contact (office):   & \dkTParidyumContactOffice{} --- \dkOrgEmail \\[0.5em]
Institution contact (office): & \dkTInstitutionContactOffice \\[0.5em]
Institution contact mailbox: & \dkTInstitutionContactMailbox
\end{tabularx}
% Both contacts are named by OFFICE, not by person (ADR-0087), and the institution's mailbox is a
% role mailbox: docket holds no person's data for an institution.

\section*{Schedule B --- Permitted Use of Names and Logos}
\textit{Anything not listed here is not permitted. ``None'' is a valid entry.}

\vspace{0.5em}
\begin{tabularx}{\linewidth}{@{}F{4.6cm}X@{}}
Paridyum may use:     & \dkTParidyumMayUse \\[0.5em]
Where and how:        & \dkTParidyumMayUseWhere \\[0.8em]
Institution may use:  & \dkTInstitutionMayUse \\[0.5em]
Where and how:        & \dkTInstitutionMayUseWhere
\end{tabularx}
% STD-0021: a logo is never used as social proof on a Paridyum page. Schedule B cannot permit
% that use even if both parties want it.

\end{document}
